% day3_seg4_ins_challenges/observability.tex
\begin{frame}
  \onslide<1->{
	\frametitle{Observability} 
	\framesubtitle{}
}
\begin{itemize}
  \onslide<2->{
 \item Structural observability is concerned with how $H$ maps measurements into states
}
  \onslide<3->{
  \item Observability means there is a linear path from the measurements to states
}
  \onslide<4->{
  \item Structural observability determined by $\Phi$ and $H$:
}
  \onslide<5->{
\item  If $H$ is $1\times n$, where $n$ is the size of the state vector, then the state $x_0$ 
is said to be \textbf{observable }from $n$ scalar measurements $\set{z_0,\ldots,z_{n-1}}$ if the block matrix
\[
\begin{bmatrix}
H\trans&\Phi\trans H\trans & \cdots & (\Phi\trans)^{n-1}H\trans
\end{bmatrix}
\]
is of rank $n$

}
  \onslide<6->{
 \item Will making measurements as governed by $H$ actually enable us to estimate all of the states 
 in the state vector? 
}
 \end{itemize}
\end{frame}

\begin{frame}
  \onslide<1->{
	\frametitle{Observability} 
	\framesubtitle{Designing Observability into the Filter}
}
\begin{itemize}
  \onslide<2->{
  \item Observability challenge for filters is often a matter of choosing right kinds of sensors
}
  \onslide<3->{
  \item Choosing the right measurement scheduling
}
  \onslide<4->{
  \item Certain kinds of divergence can be the result of poor observability
}
  \onslide<5->{
\item   Example: Let $H=\left[1\;\;0\;\; 0\;\; 0\;\; 0\right]$, 
$\Phi=
\begin{bmatrix}
0&1&0&0&0\\
0&0&1&0&0\\
0&0&0&1&0\\
0&0&0&0&1\\
1& 2&3&4&5
\end{bmatrix}$
}
  \onslide<6->{
	\item What is the rank 
	of $\left[H\trans\;\;\Phi H\trans\;\;\Phi^2 H\trans\;\; \Phi^3 H\trans\;\;\Phi^4 H\trans\right]$?
}
\end{itemize}
\end{frame}

